\section*{Chapter \ref{ch:rs:universe-symbology-and-corporate-actions} --- Universe, Symbology and Corporate Actions}

\begin{solution}{exo:rs:universe-symbology-and-corporate-actions:1}
$12/3 = 4$: the adjusted price is USD~4, so a floor on adjusted prices excludes a stock that traded at USD~12, because it later
split (after rising).
\end{solution}

\begin{solution}{exo:rs:universe-symbology-and-corporate-actions:2}
$24/0.80 - 1 = 29$, a return of $+2\,900\%$ between two different companies.
\end{solution}

\begin{solution}{exo:rs:universe-symbology-and-corporate-actions:3}
Each month 11.5\% of the capital is sold and 11.5\% bought, at 0.2\% each: $12 \times 0.115 \times 2 \times 0.002 = 0.55\%$ of capital a
year. With the buffer, $12 \times 0.031 \times 0.004 = 0.15\%$.
\end{solution}

\begin{solution}{exo:rs:universe-symbology-and-corporate-actions:4}
$0.02 \times 0.30 = 0.6$ points a year: the missing \omterm{def:rs:universe-symbology-and-corporate-actions:delisting}{delisting return} is lost on exactly the positions that fail.
\end{solution}

\begin{solution}{exo:rs:universe-symbology-and-corporate-actions:5}
Ranks near the boundary move by a few places a month in both directions; a name at rank 499 is as likely to be at 502 next month as
at 496, and without a buffer every such crossing is a replacement. With an exit rank of 600 a member must fall a hundred places to
leave, which takes a real change in its liquidity. The members kept are still in the top 600, so the universe's liquidity is almost
unchanged while the pointless replacements disappear.
\end{solution}

\begin{solution}{exo:rs:universe-symbology-and-corporate-actions:6}
The split-adjusted series: the distance is a ratio of two prices in the window, and later factors cancel in it, while splits inside
the window are removed. Traded prices turn a split inside the window into an apparent fall of 50\% or more; a total-return series
mixes dividends into a price level, so a high-dividend stock looks further below its high than its price is.
\end{solution}

\begin{solution}{exo:rs:universe-symbology-and-corporate-actions:7}
With probability 0, no reused ticker and no spliced return; with 0.35, 164 spliced returns (158 tickers reused), mean absolute size
164\% and median 71\%; with 0.7, 398 (371 tickers), mean 259\% and median 90\%. The number grows with the probability; each size is the
ratio of two unrelated prices, and the mean is dominated by the few cases where the old holder's last price was near zero (a
performance delisting), which the median ignores.
\end{solution}

\begin{solution}{exo:rs:universe-symbology-and-corporate-actions:8}
The current file contains only today's survivors (\omterm{def:rs:point-in-time-data-and-the-biases:survivorship}{survivorship bias}) and its adjusted prices carry later splits into the past (the
floor drops future winners); rebuilding back to 2010 from it applies both biases at every date. The universe must come from a
point-in-time \omterm{def:rs:universe-symbology-and-corporate-actions:master}{security master} with delisted names, the floor on traded prices, and the history requirement counted from each date
backwards.
\end{solution}

\begin{solution}{pb:rs:universe-symbology-and-corporate-actions:1}
\textbf{1.} 1\,471 listings, 471 of which end. \textbf{2.} 113 for performance, 358 by merger. \textbf{3.} 335 splits in 231 names.
\textbf{4.} 158 of 1\,549 tickers. \textbf{5.} 164. \textbf{6.} 164\% (median 71\%). \textbf{7.} $39.55/4.23 - 1 = +835\%$.
\textbf{8.} 0.30\%. \textbf{9.} The spliced windows are only those that straddle a switch, and the liquid universe rarely holds the
failing companies before their delisting or the new listings soon after theirs. \textbf{10.} 0.34\%. \textbf{11.} 71\% against
15\%. \textbf{12.} A stock splits because its price rose; the adjusted price divides the early prices by the later split, so the
stocks pushed below the floor are the ones that went on to rise. \textbf{13.} 8.65 traded and 4.32 adjusted at month 0; a two-for-one
split in month 57. \textbf{14.} Any rule on a level: a capitalisation cut-off computed with adjusted prices and current shares, a
minimum price for short sales, a distance from a past high computed with total-return prices. \textbf{15.} 11.5\% and 3.1\% a month.
\textbf{16.} Almost nothing: every member is within the top 600 by traded value. \textbf{17.} The simulation books the reason's
return on the delisting month ($-30\%$ or the merger price); in real data from the \omterm{def:rs:universe-symbology-and-corporate-actions:master}{security master}'s delisting record, with an
estimate by reason when the return is missing. \textbf{18.} \emph{Named result:} joined on tickers, the simulated market's price file
contains 164 returns spanning two companies, of 164\% in absolute size on average; a USD~5 floor on adjusted prices drops 0.34\% of
name-months whose next-year return is 71\%, against 15\% for the universe. \textbf{19.} Any table keyed on tickers, and any
``current'' constituent list used for the past. \textbf{20.} A ticker names a listing only for a date range, and the key must not
change meaning over time.
\end{solution}

\begin{solution}{iq:rs:universe-symbology-and-corporate-actions:1}
Tickers change (renames) and are reused after delistings, so one ticker can name several companies over time and one company
several tickers. Key everything on a permanent internal identifier, with a mapping from each external symbol to it by date range.
\iqlookfor{renames and reuse; a dated mapping to a permanent key.}
\end{solution}

\begin{solution}{iq:rs:universe-symbology-and-corporate-actions:2}
Look-ahead through later splits: the adjusted price at an early date is divided by the splits that followed, which happen after
prices rise, so the floor excludes future winners. Apply level rules to the price that traded that day.
\iqlookfor{that adjustment factors come from the future.}
\end{solution}

\begin{solution}{iq:rs:universe-symbology-and-corporate-actions:3}
A securities table (permanent id, issuer, listing and delisting dates, delisting reason and return) and a symbol table (symbol type,
symbol, permanent id, start date, end date) indexed on (symbol, start date); the query takes the row with start $\le D <$ end. Add a
corporate-actions table keyed on permanent id with announcement and ex-dates.
\iqlookfor{date-ranged rows and an index on (symbol, start).}
\end{solution}

\begin{solution}{iq:rs:universe-symbology-and-corporate-actions:4}
From capacity and costs: large enough for breadth (chapter~15) and small enough that every name can be traded at the strategy's
size, with liquidity measured as known at each date; a buffer (exit rank wider than the size) and a minimum history to limit churn
and cold starts; no rules on adjusted levels.
\iqlookfor{liquidity versus breadth, a buffer, point-in-time rules.}
\end{solution}

\begin{solution}{iq:rs:universe-symbology-and-corporate-actions:5}
The return to what a holder could realise: the over-the-counter price after delisting, or the recovery. If it is missing, an
estimate by delisting reason (such as the $-30\%$ Shumway recovered for performance delistings), flagged; never zero, which is the
optimistic error on the losing positions.
\iqlookfor{the \omterm{def:rs:universe-symbology-and-corporate-actions:delisting}{delisting return}, and not booking zero.}
\end{solution}

\begin{solution}{iq:rs:universe-symbology-and-corporate-actions:6}
Between actions all prices are divided by one constant, which cancels in a return; at an ex-date the factor removes the mechanical
price change, so the adjusted return equals the holding return. A floor on adjusted prices at date $t$ uses the split factors of
actions after $t$, which are more likely for stocks that later rose.
\iqlookfor{the cancellation argument and a concrete leak.}
\end{solution}
